\section*{अध्याय \ref{ch:g12:stereochemistry} --- त्रिविम रसायन: किरैलता और समावयवी}

\begin{solution}{exo:g12:stereochemistry:1}
ब्यूटेन-2-ऑल: कार्बन 2 (\ce{H}, \ce{OH}, \ce{CH3}, \ce{CH2CH3})।
प्रोपेन-2-ऑल: कोई नहीं (कार्बन 2 पर दो एक जैसे \ce{CH3} हैं)।
2-क्लोरोब्यूटेन: कार्बन 2। लैक्टिक अम्ल: कार्बन 2 (\ce{H}, \ce{OH},
\ce{CH3}, \ce{COOH})।
\end{solution}

\begin{solution}{exo:g12:stereochemistry:2}
पहला: दोनों \ce{Cl} (\ce{H} पर प्राथमिकता) एक ही ओर: \omterm{def:g12:stereochemistry:z-e}{Z}।
दूसरा: बाईं ओर \ce{CH3} को प्राथमिकता है, ऊपरी ओर; दाईं ओर
\ce{CH2CH3} को प्राथमिकता है, निचली ओर: \omterm{def:g12:stereochemistry:z-e}{E}।
\end{solution}

\begin{solution}{exo:g12:stereochemistry:3}
\omterm{def:g12:stereochemistry:chiral}{किरैल}: दस्ताना, पेंच। \omterm{def:g12:stereochemistry:chiral}{किरैल} नहीं: चम्मच, घन (हर एक में
दर्पण-तल है)। \ce{CHFClBr} \omterm{def:g12:stereochemistry:chiral}{किरैल} है (कार्बन पर चार भिन्न \omterm{def:g7:atoms-and-molecules:atom}{परमाणु});
\ce{CH2Cl2} नहीं।
\end{solution}

\begin{solution}{exo:g12:stereochemistry:4}
दो \omterm{def:g12:stereochemistry:enantiomer}{प्रतिबिंबरूपियों} की समान मात्राओं का \omterm{def:g6:pure-substances-and-mixtures:mixture}{मिश्रण}। रेसेमिक कार्वोन में
दोनों \omterm{def:g12:stereochemistry:enantiomer}{प्रतिबिंबरूपी} होते हैं, इसलिए नाक तक दोनों गंधें एक साथ पहुँचती हैं।
\end{solution}

\begin{solution}{exo:g12:stereochemistry:5}
नहीं: कार्बन 1 पर दो एक जैसे \omterm{def:g7:atoms-and-molecules:atom}{परमाणु} हैं, \ce{H} और \ce{H}; उन्हें आपस में बदलने पर
वही \omterm{def:g7:atoms-and-molecules:molecule}{अणु} मिलता है।
\end{solution}

\begin{solution}{exo:g12:stereochemistry:6}
\setchemfig{atom sep=2em}
\chemfig{C(-[2]CH_2CH_3)(-[4]H_3C)(<[:-30]OH)(<:[:-150]H)} \quad दर्पण
\quad \chemfig{C(-[2]CH_2CH_3)(-[0]CH_3)(<[:-150]HO)(<:[:-30]H)}
\end{solution}

\begin{solution}{exo:g12:stereochemistry:7}
2-क्लोरोब्यूटेन: 2 (एक \omterm{def:g12:stereochemistry:chiral}{असममित कार्बन})। 2,3-डाइक्लोरोब्यूटेन: 3 (एक जैसे अर्धों वाले
दो \omterm{def:g12:stereochemistry:chiral}{असममित कार्बन}: \omterm{def:g12:stereochemistry:enantiomer}{प्रतिबिंबरूपियों} का एक जोड़ा और एक मेसो
रूप)। पेंट-2-ईन: 2 (\omterm{def:g12:stereochemistry:z-e}{Z} और \omterm{def:g12:stereochemistry:z-e}{E})।
\end{solution}

\begin{solution}{exo:g12:stereochemistry:8}
\omterm{def:g12:stereochemistry:enantiomer}{प्रतिबिंबरूपी}; \omterm{def:g12:stereochemistry:diastereomer}{अप्रतिबिंबरूपी}; \omterm{def:g12:stereochemistry:diastereomer}{अप्रतिबिंबरूपी}; एक ही \omterm{def:g7:atoms-and-molecules:molecule}{अणु} (मेसो-टार्टरिक
अम्ल अपने दर्पण-प्रतिबिंब पर अध्यारोपित हो जाता है)।
\end{solution}

\begin{solution}{exo:g12:stereochemistry:9}
विपरीत \ce{CH3} समूहों वाले सांतरित \omterm{def:g12:stereochemistry:conformation}{संरूपण} में बड़े
समूह एक-दूसरे से यथासंभव दूर होते हैं: यह सबसे स्थायी है।
ग्रसित \omterm{def:g12:stereochemistry:conformation}{संरूपण} में दोनों \ce{CH3} आमने-सामने होते हैं और एक-दूसरे को भीड़ देते हैं।
\end{solution}

\begin{solution}{exo:g12:stereochemistry:10}
उसके दोनों अर्ध, हर एक \ce{CH(OH)COOH}, एक-दूसरे के दर्पण-प्रतिबिंब हैं:
दोनों केंद्रीय कार्बनों के बीच का एक तल \omterm{def:g7:atoms-and-molecules:molecule}{अणु} को दो दर्पण-अर्धों में
काटता है। दर्पण-तल वाला \omterm{def:g7:atoms-and-molecules:molecule}{अणु} स्वयं अपना दर्पण-प्रतिबिंब होता है: वह
\omterm{def:g12:stereochemistry:chiral}{किरैल} नहीं।
\end{solution}

\begin{solution}{exo:g12:stereochemistry:11}
\setchemfig{atom sep=1.7em}
(\omterm{def:g12:stereochemistry:z-e}{Z})-पेंट-2-ईन \chemfig{C(-[:150]H_3C)(-[:210]H)=C(-[:30]CH_2-CH_3)(-[:-30]H)}
और (\omterm{def:g12:stereochemistry:z-e}{E})-पेंट-2-ईन \chemfig{C(-[:150]H_3C)(-[:210]H)=C(-[:30]H)(-[:-30]CH_2-CH_3)}।
\end{solution}

\begin{solution}{exo:g12:stereochemistry:12}
टार्टरिक अम्ल की तरह, पर \ce{Br} ने \ce{OH} की जगह ली है और
\ce{CH3} ने \ce{COOH} की: दोनों \ce{Br} ठोस, दोनों डैश वाले (\omterm{def:g12:stereochemistry:enantiomer}{प्रतिबिंबरूपियों}
का एक जोड़ा), और एक ठोस एक डैश वाला, जिसमें दर्पण-तल है:
मेसो रूप। तीन \omterm{def:g12:stereochemistry:stereoisomer}{त्रिविम समावयवी}।
\end{solution}

\begin{solution}{exo:g12:stereochemistry:13}
आधा। रसायनज्ञ दोनों \omterm{def:g12:stereochemistry:enantiomer}{प्रतिबिंबरूपियों} को अलग कर सकता है (जैसे पाश्चर ने किया, पर
दूसरे साधनों से), या केवल उपयोगी \omterm{def:g12:stereochemistry:enantiomer}{प्रतिबिंबरूपी} बना सकता है, ऐसी अभिक्रिया से जो
स्वयं \omterm{def:g12:stereochemistry:chiral}{किरैल} हो (एक \omterm{def:g12:stereochemistry:chiral}{किरैल} उत्प्रेरक, एक एंज़ाइम, एक \omterm{def:g12:stereochemistry:chiral}{किरैल} प्रारंभिक
पदार्थ)।
\end{solution}

\begin{solution}{exo:g12:stereochemistry:14}
तीन: (2E,4E), (2Z,4Z) और (2E,4Z), जिनमें अंतिम वही \omterm{def:g7:atoms-and-molecules:molecule}{अणु} है जो
दूसरे सिरे से पढ़ा गया (2Z,4E)।
\end{solution}

\begin{solution}{exo:g12:stereochemistry:15}
$2^3 = 8$ \omterm{def:g12:stereochemistry:stereoisomer}{त्रिविम समावयवी}, अर्थात \omterm{def:g12:stereochemistry:enantiomer}{प्रतिबिंबरूपियों} के 4 जोड़े। हर
\omterm{def:g12:stereochemistry:stereoisomer}{त्रिविम समावयवी} का एक \omterm{def:g12:stereochemistry:enantiomer}{प्रतिबिंबरूपी} और $8 - 2 = 6$ \omterm{def:g12:stereochemistry:diastereomer}{अप्रतिबिंबरूपी} हैं।
\end{solution}

\begin{solution}{pb:g12:stereochemistry:1}
\textbf{1.} \ce{C4H6O6}; $M = 4 \times 12.0 + 6 \times 1.0 + 6 \times 16.0
= \qty{150.0}{g/mol}$।

\textbf{2.} दो कार्बोक्सिल समूह (अम्ल) और दो हाइड्रॉक्सिल समूह
(\omterm{def:g11:functional-groups:alcohol}{ऐल्कोहॉल})।

\textbf{3.} दोनों केंद्रीय कार्बन, हर एक \ce{H}, \ce{OH},
\ce{COOH} और \ce{CH(OH)COOH} से आबंधित।

\textbf{4.} $2^2 = 4$।

\textbf{5.} दोनों ठोस पच्चरों से।

\textbf{6.} दोनों \ce{OH} डैश वाले पच्चरों से; इसे L पर अध्यारोपित नहीं किया जा
सकता: यह D-टार्टरिक अम्ल है, उसका \omterm{def:g12:stereochemistry:enantiomer}{प्रतिबिंबरूपी}।

\textbf{7.} केंद्रीय आबंध के चारों ओर इस तरह घुमाने पर कि दोनों \ce{OH} एक ही
दिशा में हों, \omterm{def:g7:atoms-and-molecules:molecule}{अणु} में उसके दोनों केंद्रीय कार्बनों के बीच एक तल होता है
जो एक अर्ध को दूसरे पर प्रतिबिंबित करता है।

\textbf{8.} नहीं: यह मेसो-टार्टरिक अम्ल है।

\textbf{9.} चार संयोजनों (ठोस, ठोस), (डैश, डैश),
(ठोस, डैश), (डैश, ठोस) में से अंतिम दो एक ही \omterm{def:g7:atoms-and-molecules:molecule}{अणु} हैं,
मेसो रूप, क्योंकि \omterm{def:g7:atoms-and-molecules:molecule}{अणु} के दोनों अर्ध एक जैसे हैं।

\textbf{10.} L और D: \omterm{def:g12:stereochemistry:enantiomer}{प्रतिबिंबरूपी}। L और मेसो: \omterm{def:g12:stereochemistry:diastereomer}{अप्रतिबिंबरूपी}।

\textbf{11.} \qty{169}{\celsius} पर ही: \omterm{def:g12:stereochemistry:enantiomer}{प्रतिबिंबरूपियों} के
भौतिक गुण समान होते हैं।

\textbf{12.} नहीं: यह L का \omterm{def:g12:stereochemistry:diastereomer}{अप्रतिबिंबरूपी} है, अपने गुणों वाला एक
भिन्न यौगिक।

\textbf{13.} नाक के ग्राही \omterm{def:g12:stereochemistry:chiral}{किरैल} होते हैं और दर्पण-प्रतिबिंबों में
भेद करते हैं; थर्मामीटर, जो \omterm{def:g12:stereochemistry:chiral}{किरैल} नहीं, ऐसा नहीं कर सकता।

\textbf{14.} नहीं, यह दो यौगिकों का \omterm{def:g6:pure-substances-and-mixtures:mixture}{मिश्रण} है; समग्र रूप से यह
ध्रुवण-घूर्णक नहीं है, क्योंकि दोनों \omterm{def:g12:stereochemistry:enantiomer}{प्रतिबिंबरूपी} एक-दूसरे का प्रभाव निरस्त कर देते हैं।

\textbf{15.} केवल एक \omterm{def:g12:stereochemistry:enantiomer}{प्रतिबिंबरूपी}: हर क्रिस्टल केवल L से या केवल D से
बना था।

\textbf{16.} हर एक का \qty{5.0}{g}।

\textbf{17.} मेसो रूप एक भिन्न यौगिक है, L और D के \omterm{def:g12:stereochemistry:enantiomer}{रेसेमिक
मिश्रण} का भाग नहीं।

\textbf{18.} हर ढेर को घोलकर: विलयन ध्रुवित प्रकाश को
समान मात्रा में विपरीत दिशाओं में घुमाते हैं (और क्रिस्टल एक-दूसरे के दर्पण-प्रतिबिंब
होते हैं)।

\textbf{19.} चार: उसके दोनों \omterm{def:g12:stereochemistry:chiral}{असममित कार्बनों} पर भिन्न समूह हैं
(एक पर \ce{Cl}, दूसरे पर \ce{OH}), इसलिए किसी भी संयोजन में दर्पण-तल
नहीं है, और कोई भी दो संयोजन एक ही \omterm{def:g7:atoms-and-molecules:molecule}{अणु} नहीं हैं।

\textbf{20.} तीन: L, D और मेसो।
\end{solution}
